% !TeX root = ../../Book3.tex
% 纲要标题：### 19.6*　曲率能量与 Willmore 曲面
% 纲要位置：几何学计划纲要.md，第 2908 行。

% 纲要内容开始
% 
% > 写作范围与外部输入：规划任务：以欧氏三空间中的闭光滑浸入曲面为主范围，明确 W 为平均曲率平方的面积积分，并与第一卷第8章的法向、平均曲率和非负 Laplace 约定逐项核对。完整推导法向一阶变分与临界方程，计算球面，并用 Gauss–Bonnet 证明球面拓扑下的能量下界及等号。一般拓扑的下界、弱紧性、极小化存在、Willmore 流与环面定理的深层证明外引，不把球面模型扩称一般结论。
% 
% * 闭曲面的曲率能量与缩放
% * 法向变分与 Willmore 方程
% * 球面模型与球面拓扑下的能量下界
% * 等号条件与全脐曲面
% * 一般拓扑、弱极限与存在性的外部输入
% 
% ---
% 
% 纲要内容结束
% 正文状态：未撰写；本文件目前是框架。
% 规划任务：以欧氏三空间中的闭光滑浸入曲面为主范围，明确 W 为平均曲率平方的面积积分，并与第一卷第8章的法向、平均曲率和非负 Laplace 约定逐项核对。完整推导法向一阶变分与临界方程，计算球面，并用 Gauss–Bonnet 证明球面拓扑下的能量下界及等号。一般拓扑的下界、弱紧性、极小化存在、Willmore 流与环面定理的深层证明外引，不把球面模型扩称一般结论。

\OptionalSection{曲率能量与 Willmore 曲面}{g3:sec:1906}

待撰写

\subsection{闭曲面的曲率能量与缩放}
\label{g3:subsec:1906:01}
待撰写

\subsection{法向变分与 Willmore 方程}
\label{g3:subsec:1906:02}
待撰写

\subsection{球面模型与球面拓扑下的能量下界}
\label{g3:subsec:1906:03}
待撰写

\subsection{等号条件与全脐曲面}
\label{g3:subsec:1906:04}
待撰写

\subsection{一般拓扑、弱极限与存在性的外部输入}
\label{g3:subsec:1906:05}
待撰写
